\section{Banco de questões — Segurança da Informação}

\subsection*{N1 — Fundamentos}
\begin{questao}{\nivel{N1} 13.01}Confidencialidade busca garantir que:\begin{enumerate}[label=\Alph*)]\item informação seja acessada apenas por entidades autorizadas\item dados nunca sejam apagados\item sistemas nunca falhem\item logs sejam públicos\item toda informação seja anônima\end{enumerate}\end{questao}
\begin{questao}{\nivel{N1} 13.02}Integridade relaciona-se principalmente a:\begin{enumerate}[label=\Alph*)]\item proteção contra alteração/destruição não autorizada e preservação de exatidão\item tempo de resposta\item sigilo exclusivamente\item capacidade de armazenamento\item redundância elétrica\end{enumerate}\end{questao}
\begin{questao}{\nivel{N1} 13.03}Qual alternativa representa autenticação multifator?\begin{enumerate}[label=\Alph*)]\item senha + token físico\item senha + PIN\item duas senhas\item duas perguntas secretas\item PIN + palavra-chave\end{enumerate}\end{questao}
\begin{questao}{\nivel{N1} 13.04}Malware capaz de se propagar autonomamente pela rede é tipicamente:\begin{enumerate}[label=\Alph*)]\item worm\item trojan\item hoax\item adware apenas\item certificado\end{enumerate}\end{questao}
\begin{questao}{\nivel{N1} 13.05}Phishing é:\begin{enumerate}[label=\Alph*)]\item engenharia social para induzir vítima a revelar dados ou executar ação\item RAID com falha\item algoritmo de hash\item protocolo de roteamento\item política de backup\end{enumerate}\end{questao}
\begin{questao}{\nivel{N1} 13.06}Um IDS tem como função primária:\begin{enumerate}[label=\Alph*)]\item detectar atividade suspeita e gerar alerta\item cifrar discos\item autenticar usuários\item distribuir IP\item realizar backup\end{enumerate}\end{questao}
\begin{questao}{\nivel{N1} 13.07}Criptografia simétrica utiliza:\begin{enumerate}[label=\Alph*)]\item segredo compartilhado para cifrar/decifrar\item par público/privado obrigatoriamente\item somente hash\item certificado sem chave\item DNSSEC apenas\end{enumerate}\end{questao}
\begin{questao}{\nivel{N1} 13.08}Um hash criptográfico é, em regra:\begin{enumerate}[label=\Alph*)]\item função unidirecional que produz resumo de tamanho definido\item cifra reversível com mesma chave\item protocolo de VPN\item mecanismo de backup\item assinatura por si só\end{enumerate}\end{questao}
\begin{questao}{\nivel{N1} 13.09}PKI tem como objetivo central:\begin{enumerate}[label=\Alph*)]\item gerenciar confiança, certificados e associação entre identidades e chaves públicas\item gerenciar VLAN\item agendar processos\item balancear tráfego\item normalizar bancos\end{enumerate}\end{questao}
\begin{questao}{\nivel{N1} 13.10}RPO de 30 minutos significa, de forma simplificada:\begin{enumerate}[label=\Alph*)]\item perda temporal de dados tolerável de até cerca de 30 minutos\item retorno do serviço em 30 minutos\item duração do backup\item senha válida por 30 minutos\item sessão VPN de 30 minutos\end{enumerate}\end{questao}

\subsection*{N2 — Aplicação}
\begin{questao}{\nivel{N2} 13.11}Uma pessoa usa senha e PIN. Isso constitui MFA?\begin{enumerate}[label=\Alph*)]\item não necessariamente, pois ambos são fator de conhecimento\item sim, sempre\item sim, se tiverem 8 caracteres\item não, pois MFA só usa biometria\item sim, se houver TLS\end{enumerate}\end{questao}
\begin{questao}{\nivel{N2} 13.12}RBAC concede permissões principalmente com base em:\begin{enumerate}[label=\Alph*)]\item papéis atribuídos aos usuários\item endereço MAC\item temperatura do servidor\item número de série\item hash do arquivo\end{enumerate}\end{questao}
\begin{questao}{\nivel{N2} 13.13}Privilégio mínimo recomenda:\begin{enumerate}[label=\Alph*)]\item conceder somente acessos necessários à função e pelo tempo adequado\item dar administrador a todos\item compartilhar conta\item impedir logs\item usar uma única senha institucional\end{enumerate}\end{questao}
\begin{questao}{\nivel{N2} 13.14}Credential stuffing utiliza principalmente:\begin{enumerate}[label=\Alph*)]\item pares usuário/senha vazados de outros serviços\item exploit de buffer overflow necessariamente\item certificado revogado\item DNSSEC\item RAID\end{enumerate}\end{questao}
\begin{questao}{\nivel{N2} 13.15}SQL injection é mitigada diretamente por:\begin{enumerate}[label=\Alph*)]\item consultas parametrizadas/prepared statements e privilégio mínimo\item concatenar entrada do usuário\item exibir erro SQL completo\item executar banco como root\item desabilitar validação\end{enumerate}\end{questao}
\begin{questao}{\nivel{N2} 13.16}DMZ é usada para:\begin{enumerate}[label=\Alph*)]\item separar serviços expostos da rede interna por controles de rede\item eliminar firewall\item armazenar backups offline necessariamente\item criar hashes\item substituir autenticação\end{enumerate}\end{questao}
\begin{questao}{\nivel{N2} 13.17}Qual controle opera inline e pode bloquear tráfego malicioso conforme política?\begin{enumerate}[label=\Alph*)]\item IPS\item IDS passivo\item SIEM apenas\item honeypot\item hash\end{enumerate}\end{questao}
\begin{questao}{\nivel{N2} 13.18}Uma VPN protege principalmente:\begin{enumerate}[label=\Alph*)]\item comunicação sobre rede não confiável por túnel/políticas criptográficas conforme tecnologia\item integridade do endpoint comprometido automaticamente\item disponibilidade física do link\item senha do usuário sem autenticação\item backup do sistema\end{enumerate}\end{questao}
\begin{questao}{\nivel{N2} 13.19}Assinatura digital fornece diretamente:\begin{enumerate}[label=\Alph*)]\item autenticidade e integridade vinculadas à chave do signatário\item confidencialidade obrigatória do conteúdo\item disponibilidade\item compressão\item anonimato\end{enumerate}\end{questao}
\begin{questao}{\nivel{N2} 13.20}AES é algoritmo:\begin{enumerate}[label=\Alph*)]\item simétrico\item assimétrico\item hash\item assinatura exclusivamente\item compressão\end{enumerate}\end{questao}
\begin{questao}{\nivel{N2} 13.21}RSA é associado a criptografia:\begin{enumerate}[label=\Alph*)]\item assimétrica\item simétrica por bloco exclusivamente\item hash\item RAID\item VPN sem chaves\end{enumerate}\end{questao}
\begin{questao}{\nivel{N2} 13.22}CRL e OCSP estão relacionados a:\begin{enumerate}[label=\Alph*)]\item status/revogação de certificados\item VLAN\item controle de processos\item backup incremental\item DHCP\end{enumerate}\end{questao}
\begin{questao}{\nivel{N2} 13.23}Um honeypot deve ser:\begin{enumerate}[label=\Alph*)]\item monitorado e isolado para não virar plataforma de ataque\item colocado com acesso irrestrito à produção\item usado como backup\item substituto do firewall\item desconhecido pelos administradores\end{enumerate}\end{questao}
\begin{questao}{\nivel{N2} 13.24}Risco residual é:\begin{enumerate}[label=\Alph*)]\item risco que permanece após aplicação dos controles\item risco antes de controles\item custo do ativo\item somente ameaça externa\item vulnerabilidade sem impacto\end{enumerate}\end{questao}
\begin{questao}{\nivel{N2} 13.25}A BIA serve principalmente para:\begin{enumerate}[label=\Alph*)]\item identificar processos críticos, dependências e impactos da interrupção ao longo do tempo\item fazer pentest\item gerar certificado\item criar VLAN\item substituir plano de continuidade\end{enumerate}\end{questao}

\subsection*{N3 — Nível de concurso}
\begin{questao}{\nivel{N3} 13.26}Uma conta administrativa é usada por cinco pessoas e os logs registram apenas o nome da conta. O problema central é:\begin{enumerate}[label=\Alph*)]\item baixa accountability e atribuição individual\item falta de disponibilidade\item excesso de criptografia\item baixa largura de banda\item ausência de RAID\end{enumerate}\end{questao}
\begin{questao}{\nivel{N3} 13.27}Um atacante altera ARP para receber tráfego destinado ao gateway. O ataque é exemplo de:\begin{enumerate}[label=\Alph*)]\item spoofing/poisoning de ARP com possibilidade de MITM\item SQL injection\item buffer overflow\item hoax\item DDoS volumétrico necessariamente\end{enumerate}\end{questao}
\begin{questao}{\nivel{N3} 13.28}Um site legítimo é substituído visualmente após exploração e páginas são modificadas. Isso é:\begin{enumerate}[label=\Alph*)]\item defacement\item pharming apenas\item sniffing\item hashing\item backup\end{enumerate}\end{questao}
\begin{questao}{\nivel{N3} 13.29}Buffer overflow pode permitir corrupção de memória e controle de fluxo. Um mitigador é:\begin{enumerate}[label=\Alph*)]\item ASLR/DEP, canários e bounds checking conforme plataforma\item concatenar strings sem limite\item desabilitar atualizações\item executar como root\item remover validação\end{enumerate}\end{questao}
\begin{questao}{\nivel{N3} 13.30}Um WAF bloqueia payload conhecido, mas a aplicação continua vulnerável a SQL injection. A conclusão correta é:\begin{enumerate}[label=\Alph*)]\item WAF é controle compensatório e não substitui correção no código\item WAF remove a vulnerabilidade do código\item banco não precisa de privilégio mínimo\item testes podem ser eliminados\item TLS corrige SQL injection\end{enumerate}\end{questao}
\begin{questao}{\nivel{N3} 13.31}Um certificado TLS expirado causa indisponibilidade mesmo com chave privada intacta. Isso mostra que:\begin{enumerate}[label=\Alph*)]\item gestão do ciclo de vida de certificados é controle operacional essencial\item PKI elimina manutenção\item validade não importa\item certificado não possui período\item TLS não usa certificados\end{enumerate}\end{questao}
\begin{questao}{\nivel{N3} 13.32}Criptografia em repouso com chave armazenada no mesmo servidor e acessível ao mesmo administrador reduz principalmente qual risco, mas não todos?\begin{enumerate}[label=\Alph*)]\item perda física de mídia pode ser mitigada, mas comprometimento do servidor/admin pode acessar chave e dados\item elimina insider\item elimina ransomware\item elimina necessidade de IAM\item garante anonimização\end{enumerate}\end{questao}
\begin{questao}{\nivel{N3} 13.33}Assinar o hash de um arquivo ajuda a detectar alteração porque:\begin{enumerate}[label=\Alph*)]\item qualquer mudança relevante altera o digest e invalida verificação da assinatura\item hash é reversível\item assinatura cifra o arquivo sempre\item certificado substitui hash\item assinatura aumenta disponibilidade\end{enumerate}\end{questao}
\begin{questao}{\nivel{N3} 13.34}Uma organização implementa backup online acessível pelas mesmas credenciais de domínio usadas na produção. Em ransomware com comprometimento do domínio, o risco é:\begin{enumerate}[label=\Alph*)]\item atacante alcançar e destruir/criptografar também backups\item backup torna-se offline automaticamente\item RPO fica zero\item RAID bloqueia ataque\item MFA no usuário final elimina risco\end{enumerate}\end{questao}
\begin{questao}{\nivel{N3} 13.35}Um RTO de 1 hora e restore test de 6 horas indicam:\begin{enumerate}[label=\Alph*)]\item incompatibilidade entre capacidade demonstrada e objetivo de recuperação\item RTO cumprido\item RPO de 6 horas\item ausência de risco\item problema apenas de confidencialidade\end{enumerate}\end{questao}
\begin{questao}{\nivel{N3} 13.36}Na gestão de risco, aceitar um risco significa:\begin{enumerate}[label=\Alph*)]\item reconhecer o risco residual dentro de critérios/apetite e assumir sua consequência com decisão autorizada\item ignorá-lo sem análise\item transferi-lo automaticamente\item eliminar vulnerabilidade\item garantir que não ocorrerá\end{enumerate}\end{questao}
\begin{questao}{\nivel{N3} 13.37}ISO/IEC 27001 diferencia-se da ISO/IEC 27002 porque:\begin{enumerate}[label=\Alph*)]\item 27001 contém requisitos do SGSI; 27002 fornece orientação de controles\item ambas são idênticas\item 27002 é requisito de certificação do SGSI\item 27001 é só criptografia\item 27002 é norma de roteamento\end{enumerate}\end{questao}
\begin{questao}{\nivel{N3} 13.38}ISO/IEC 27005 trata principalmente de:\begin{enumerate}[label=\Alph*)]\item gestão de riscos de segurança da informação\item gerenciamento de serviços\item desenvolvimento ágil\item data warehouse\item cabeamento\end{enumerate}\end{questao}
\begin{questao}{\nivel{N3} 13.39}Um SIEM recebe logs, mas relógios dos sistemas variam 20 minutos. O impacto é:\begin{enumerate}[label=\Alph*)]\item correlação temporal e investigação ficam comprometidas\item aumenta confidencialidade\item reduz necessidade de retenção\item impede hash\item melhora detecção\end{enumerate}\end{questao}
\begin{questao}{\nivel{N3} 13.40}Uma política de senha exige troca a cada 15 dias, levando usuários a variações previsíveis. A avaliação madura é:\begin{enumerate}[label=\Alph*)]\item controle pode gerar comportamento fraco; política deve considerar risco, MFA, credenciais comprometidas e orientação atual\item trocas frequentes sempre aumentam segurança\item MFA torna senha irrelevante\item senha nunca precisa de gestão\item complexidade é único fator\end{enumerate}\end{questao}

\subsection*{N4 — Avançado / Auditoria}
\begin{questao}{\nivel{N4} 13.41}Auditoria encontra administradores capazes de alterar registros de negócio e apagar seus próprios logs. A prioridade é:\begin{enumerate}[label=\Alph*)]\item segregar funções e proteger trilha contra alteração pelo usuário monitorado\item adicionar RAID 6\item aumentar largura de banda\item usar mais VLANs sem política\item remover logs\end{enumerate}\end{questao}
\begin{questao}{\nivel{N4} 13.42}Após ransomware, backups foram criptografados porque estavam montados permanentemente. A recomendação mais completa é:\begin{enumerate}[label=\Alph*)]\item cópias offline/imutáveis, segregação de credenciais, teste de restore e monitoramento de backup\item apenas aumentar espaço\item manter backup montado e trocar extensão\item usar RAID 0\item eliminar versionamento\end{enumerate}\end{questao}
\begin{questao}{\nivel{N4} 13.43}Um órgão usa MFA por push e registra várias aprovações fraudulentas após bombardeio de notificações. Qual controle adicional é adequado?\begin{enumerate}[label=\Alph*)]\item number matching/phishing-resistant MFA, limitação de prompts e detecção de MFA fatigue\item aumentar quantidade de pushes\item remover MFA\item usar senha compartilhada\item desabilitar alertas\end{enumerate}\end{questao}
\begin{questao}{\nivel{N4} 13.44}Uma VPN permite acesso remoto à rede inteira após autenticação, sem verificar postura do dispositivo. O risco principal é:\begin{enumerate}[label=\Alph*)]\item endpoint comprometido ganha amplo movimento lateral; segmentação e acesso contextual devem ser avaliados\item VPN garante integridade do endpoint\item criptografia elimina malware\item não há risco interno\item DNS impede movimento lateral\end{enumerate}\end{questao}
\begin{questao}{\nivel{N4} 13.45}Um fornecedor afirma conformidade ISO 27001 porque o provedor de nuvem possui certificado. A conclusão correta é:\begin{enumerate}[label=\Alph*)]\item certificação do provedor não comprova que o sistema/cliente esteja em conformidade; escopo e responsabilidades devem ser avaliados\item toda aplicação do cliente fica certificada\item conformidade é herdada integralmente\item controles do cliente são desnecessários\item certificado substitui auditoria\end{enumerate}\end{questao}
\begin{questao}{\nivel{N4} 13.46}A matriz de risco classifica ransomware como baixo impacto porque só considera custo de hardware. Qual falha metodológica?\begin{enumerate}[label=\Alph*)]\item impacto ignora indisponibilidade, dados, reputação, sanções e missão institucional\item risco deve considerar só hardware\item probabilidade substitui impacto\item confidencialidade não importa\item ransomware não afeta serviço\end{enumerate}\end{questao}
\begin{questao}{\nivel{N4} 13.47}Um restore test restaura arquivos, mas não valida aplicação, dependências nem consistência do banco. O auditor deve concluir que:\begin{enumerate}[label=\Alph*)]\item evidência é parcial; recuperação do serviço precisa testar cadeia completa e RTO/RPO\item DR está plenamente validado\item arquivo restaurado prova qualquer aplicação\item RTO é irrelevante\item backup substitui continuidade\end{enumerate}\end{questao}
\begin{questao}{\nivel{N4} 13.48}Logs de autenticação mostram milhares de tentativas distribuídas, cada IP tentando poucas senhas contra muitas contas. Isso é compatível com:\begin{enumerate}[label=\Alph*)]\item password spraying distribuído, que evita limiar por conta/IP simples\item SQL injection\item buffer overflow\item RAID rebuild\item DNSSEC\end{enumerate}\end{questao}
\begin{questao}{\nivel{N4} 13.49}Um honeypot comprometido possui rota irrestrita para a Internet e redes internas. Qual problema?\begin{enumerate}[label=\Alph*)]\item o ativo de engodo pode ser usado como pivô; isolamento e egress control são necessários\item honeypot deve ter privilégio máximo\item é impossível comprometer honeypot\item monitoramento é desnecessário\item honeypot substitui IPS\end{enumerate}\end{questao}
\begin{questao}{\nivel{N4} 13.50}Em investigação, um arquivo possui hash preservado desde a coleta, mas a origem do arquivo não foi autenticada. O que o hash prova?\begin{enumerate}[label=\Alph*)]\item que a cópia não mudou desde o hash, não que o conteúdo original era verdadeiro\item que a fonte era legítima\item autoria do arquivo\item confidencialidade\item ausência de malware\end{enumerate}\end{questao}
